En este capítulo, nos enfocaremos en establecer los pilares fundamentales de nuestra investigación, delineando tanto el objetivo general como los objetivos específicos que guiarán nuestro estudio. Nuestra investigación se centra en explorar las profundidades de la predicción de fallas de software, un campo que continúa evolucionando con el avance de las tecnologías de programación y análisis de datos
.
\subsection*{Objetivo General}

Determinar la efectividad de la combinación de técnicas de selección de características y estrategias de balanceo de clases en la mejora de la exactitud de los modelos de predicción de fallas de software utilizando el conjunto de datos BugHunter como caso de estudio.

\subsection*{Objetivos especificos}

A partir de este objetivo amplio, se derivan varios objetivos específicos que se centran en aspectos críticos de nuestra investigación, tales como la identificación de métricas predictivas claves, la utilizacion de algoritmos de análisis de selección de características y la evaluación de diferentes estrategias de balanceo de clases. Este enfoque nos permitirá no solo avanzar en la comprensión teórica de la predicción de fallas de software sino también ofrecer insights prácticos que puedan ser aplicados en el desarrollo y mantenimiento de software. A continuación, listamos los objetivos específicos:
\begin{enumerate}
    \item Evaluar el impacto del análisis de relevancia y control de redundancia en la mejora de la identificación de errores en el software utilizando métricas de código, con el fin de establecer cómo estas técnicas pueden  la predicción de bugs.

analizar el final del anterior
    
    \item Identificar cuáles métricas de código fuente actúan como mejores predictoras de errores, a nivel archivo, clase, método y proyecto.

    buscar analisis a nivel de proyecto en las metricas
    
    \item Analizar cómo la variación en los parámetros de entrada a los algoritmos de análisis de relevancia y control de redundancia afecta la efectividad y precisión de los modelos de clasificación.
    
    \item Identificar la estrategia de balanceo de clases más efectiva, determinando la influencia de los algoritmos de balanceo de clases Random Over-Sampling (ROS), Random Under-Sampling (RUS) y Synthetic Minority Over-sampling Technique (SMOTE) en la mejora de los resultados de  los clasificadores.
    \item Determinar las diferencias significativas y patrones comunes en la eficacia de las técnicas de predicción de fallas entre el proyecto resumen y los proyectos individuales.
\end{enumerate}
